\section{Conclusiones}
\label{sec:conclusion}

A lo largo de este trabajo práctico, se comprobó la eficacia de la red de Hopfield como modelo de memoria autoasociativa, capaz de recuperar información original a partir de entradas ruidosas o incompletas al guiar el sistema hacia los atractores aprendidos.

Sin embargo, también se evidenciaron sus limitaciones fundamentales. Por un lado, se demostró que la correlación entre los patrones aumenta la interferencia, deformando el paisaje de energía y reduciendo drásticamente la capacidad máxima de almacenamiento. Por otro lado, la pérdida de conexiones sinápticas (daño estructural), si bien es tolerada hasta cierto punto, termina aplanando los valles de energía de las memorias originales. Esto disminuye la fuerza de atracción de los patrones correctos, provocando que la red decante inevitablemente en estados espurios ante el ruido.

En síntesis, aunque el modelo ilustra de forma elegante cómo la memoria asociativa surge de reglas de conexión locales, su desempeño depende críticamente de la descorrelación de los datos enseñados y de la integridad de su matriz de pesos.